\section{第一段创业：从微信工具到商业化教训}

上一章建立路线，本章先看肖弘的第一段创业。他 2015 年从华中科技大学毕业后直接创业，早期做过校园产品但反馈不佳，后来转向微信公众号编辑器，再到企业微信 SCRM。这个阶段最大的训练，不只是做出了产品，而是从失败、现金流、商业化、资本和出售公司中学会：工具产品不是故事越大越好，而是要找到真实用户、真实付费和可持续单元经济。

\begin{figure}[H]
\centering
\includegraphics[width=0.92\textwidth]{founder-cycle.png}
\caption{连续创业者循环：技能、资本、市场、产品和心态不断出现新维度。自制概念图，依据 00:05:15--00:58:15 对谈内容整理。}
\end{figure}

\begin{knowledgebox}{读图：创业能力不是一次性点满}
图中 Founder 周围是技能点、资本、市场、产品、心态和博弈。肖弘反复说，每次觉得自己技能点还可以时，都会冒出新维度。第一段创业先学产品，后来学商业化，再后来发现资本也是新维度。
\end{knowledgebox}

\subsection{从学校正反馈到社会负反馈}

本节解释早期挫折。肖弘在大学时有技术社团、新媒体中心、博客和产品实践，获得了很多正反馈；毕业后做匿名社交、二手市场等产品却没有用户，团队也接近要去北京工作。这段经历说明，校园场景里的小成功不能直接外推到真实市场。真实创业更残酷，因为用户不是熟人圈，渠道不是天然存在，付费也不是自动发生。

\begin{warningbox}{常见误区：把校园正反馈当市场验证}
在校园里做出工具、公众号或小产品，能证明个人能力和兴趣，但不能证明商业市场成立。创业后的第一课，往往是发现自己原来获得的是局部正反馈，而不是广泛需求。
\end{warningbox}

\subsection{商业化为什么要更早开始}

本节看公众号编辑器的商业化。肖弘复盘说，团队花了太久做用户增长和产品迭代，商业化思考不足。后来他们选择像 SaaS 一样卖软件，而不是做广告撮合平台。广告撮合听起来天花板更高，但它不一定属于软件团队的核心能力；卖软件看起来更小，却更本分、更闭环，也能让团队持续改进产品。

\begin{importantbox}{实践经验：天花板高不等于能做成}
很多创业团队为了讲更大的故事，会选择不匹配自己能力的商业模式。公众号编辑器案例里，“工具获客后做广告撮合”听起来更大，但真正成立的是订阅软件。对工具产品来说，用户直接为软件付费，反而更能形成健康闭环。
\end{importantbox}

\subsection{VC 是昂贵的集资手段}

本节进入资本。肖弘说，VC 的贵不体现在公司不好的时候，而体现在公司好的时候：如果公司大成，早期股权融资的成本会被放大。融资是合理工具，但创始人必须知道它会改变故事、目标和压力。融了太多钱，可能逼迫公司讲更大的故事，做不一定正确的事情。

\begin{figure}[H]
\centering
\includegraphics[width=0.90\textwidth]{vc-expensive-capital.png}
\caption{VC 是昂贵的集资手段：贵不在困难时，而在公司变好时。自制概念图，依据 00:33:16--00:42:10 对谈内容整理。}
\end{figure}

\begin{knowledgebox}{读图：资本是工具，不是方向}
左侧困难时，融资像救命钱；右侧变好时，早期稀释和更大承诺会被放大。读这张图时要抓住：VC 不坏，关键是 Founder 要知道自己用这个工具换来了什么约束。
\end{knowledgebox}

\subsection{投资人不是老板}

本节保留一个高价值 founder voice。肖弘建议，不要把投资人当老板；投资人转一篇文章，权重应接近一个 KOL 的观点，而不是自动变成公司战略。如果投资人真的有强观点，就应该开严肃会议，记录谁持有什么判断，而不是让 Founder 猜投资人是不是想让公司转向。原因是 Founder 的杠杆率极高，一个误判会牵动资金和团队方向。

\begin{importantbox}{课堂提示：高杠杆决策要有严肃机制}
Founder 可以听投资人意见，但不能把随手建议当命令。重大方向要进入正式讨论，明确观点、证据、责任和决策机制。否则最危险的不是投资人错，而是 Founder 误读了投资人的随口表达。
\end{importantbox}

\subsection{本章小结}

第一段创业给肖弘留下三条底层经验：工具产品要早做商业化；融资是昂贵工具，不是战略本身；Founder 必须保持判断主权。这些经验后来直接影响 Monica、出海和 Manus 的判断。

